\section{Exploratory analyses: procedures}
\label{app:exploratory}

% Moved from Section 6.7 of the main body on 24 October 2026. The text of the procedures is unchanged; the tables are
% added after the main computation (24 October 2026); they are generated by make_tables.py.

% The record of the main computation (notes.md) confirms that no value of matching or utility had been opened before
% the procedures of the last three analyses were committed; the main results were first written on 24 October 2026.
% Analyses 6 and 7 were decided on 29 October 2026, after the main results; their
% procedures are in notes.md ("procedures, part 3"), committed that day, and they ran from 12 to 17 November 2026.
Seven analyses were added after the preregistration was fixed and are not part of it. They change no registered value
or procedure, are reported as exploratory and are not used to answer the four questions. The procedures of the first
five were fixed and committed to the study repository before any value of matching or utility on the human data had
been seen, the first two on 13 October 2026 and the other three on 21 October 2026, as the record of that repository
shows. The third to fifth use the strongest
attacker selected in the main computation for the same form and $N$, without selecting again, on full texts, the main
forms, $N = 10$ and $1{,}000$ and the 500 writers, with the same resamples. The last two were decided after the main
results had been seen, because the full-text accuracy on the coarse forms was far below that on windows and with the
lighter penalty (\S\ref{sec:results:forms}); their procedures were committed on 29 October 2026, before they were
run. They aim to strengthen the attackers, not to protect a form.

\paragraph{Uniqueness of word-length sequences among $N$ candidates.}
The first analysis uses the texts alone. For each of the 500 writers and each $k$ from 1 to 50, five runs of $k$
consecutive words are drawn from the writer's text, and we report the share for which none of the $N - 1$ decoys,
drawn as in the main computation, contains the same sequence of word lengths at any position; as a secondary quantity,
the accuracy $1/m$ of an attacker who chooses uniformly among the $m$ texts that contain it. Next to these curves we
draw reference lines $k^{*} = \log_2 N / h$, where $h$ is a plug-in estimate of the entropy rate of word lengths in
bits per word, conditioned on up to two preceding words~\cite{kalimeri2012entropy}; plug-in estimates fall as the
order rises, and the lines are a reference, not a prediction. We call the quantity uniqueness among $N$ candidates and
not a risk of re-identification. Whether uniqueness in a sample says anything about a population is
disputed~\cite{sanchez2016comment,rocher2019estimating}, and in location data unicity decreases, slowly, as the dataset
grows~\cite{farzanehfar2021risk}; here the population is the set of $N$ candidates itself, and we do not extend the
quantity beyond the values of $N$ tried.

\paragraph{A form with differentially private noise.}
The second analysis adds a form R-DP($\varepsilon$), $\varepsilon = 0.5$, $1$, $2$, $4$, built on R3, with letters and
digits released as one category. Two mechanisms use the same $\varepsilon$. First, the number of counted letter and
digit presses of each segment receives two-sided geometric noise with sensitivity 1, and letter presses are added at
random positions in the segment or removed at random accordingly, with the count kept at zero or above. Second, every
press interval is clipped at 2,000\,ms, the threshold of the P-bursts, and receives two-sided geometric noise with
sensitivity 2,000\,ms, and negative intervals are set to zero. Each guarantee holds for its own unit, one segment count
or one interval, and not for a writer's whole log: over a log the budgets add up over all segments and
intervals~\cite{kellaris2014event}. We do not claim that the combined release is differentially private for neighbouring
logs that differ by an inserted or removed press with its time. The noise is drawn from a seeded generator so that the
evaluation can be repeated, which is acceptable for evaluation but not for an actual release~\cite{holohan2023random},
and an actual release would also need an implementation that is exact on finite
computers~\cite{canonne2020discrete}. The attackers are those of R3, with the timing table of A3 built from training
logs with the same interval noise; the candidates are unchanged, and every process measure that R3 computes is counted as
approximated. The scope, samples and resamples are those of the main computation, with the contrast R3~$-$~R-DP and the
open-world area, but no negative control. R-DP is reported in exploratory tables, where $\varepsilon$ is labelled as the budget per segment count and per
interval, and it is not part of the main table or figure.

\paragraph{The upper tail of accuracy per writer.}
The third analysis reports the 95th and 99th percentiles of accuracy per writer and the share of writers whose true text
is chosen alone in every draw. No value per writer is reported.

\paragraph{Matches above a margin.}
The fourth analysis takes, for each log and draw, the difference between the attacker's second-lowest and lowest cost
over the candidates (zero if the lowest is tied), with log-likelihoods turned into costs by a change of sign. At
thresholds set at the 0th, 10th, \ldots, 90th percentiles of these differences, it reports the share of logs and draws
at or above the threshold and the accuracy among them. Costs of different attackers are not compared.

\paragraph{Accuracy by score band.}
The fifth analysis divides the 125 sampled writers of each prompt into thirds by their human score and reports the strongest
attacker's accuracy in each third and the difference between the top and bottom thirds. The attackers do not use the
scores.

\paragraph{A grid of the penalty.}
The sixth analysis runs A2 on R3, R4($w$) and R7($p$), and A1$'$ on R3 and R7($p$), the forms on which the main
computation runs it, with $\lambda = 2$, $4$, $8$ and $16$ in the alignment of \S\ref{sec:attackers:alignment}, which is
otherwise unchanged. The main computation left R7 out of the registered secondary analysis with $\lambda = 4$ and 16 (a deviation,
Appendix~\ref{app:checks:deviations}), so the values of R7 with $\lambda = 4$ and 16 are those of that analysis,
computed here after the main results, and those with $\lambda = 2$ are new. As positive controls, A2 and A1$'$ with $\lambda = 8$ equal the cells of the main
computation at both $N$, and A2 with $\lambda = 4$ and $16$ equals the secondary analysis on R3 and R4($w$) at
$N = 1{,}000$.

\paragraph{A full-text attacker that aligns chunks.}
The seventh analysis, A2-scan, cuts the record values of A2 into consecutive chunks of 20 segments from the start of
the log and drops a last chunk of fewer than five segments; a log with no chunk leaves all candidates tied. Each chunk
is aligned with the word lengths of a candidate, starting and ending anywhere in the text, with the groups and the
penalty of A2. A candidate's cost is the sum of the costs of the chunks, infinite if a chunk has no path, and the
candidate with the lowest cost is chosen. For simplicity, chunks may overlap or change order within a candidate. It
runs with $\lambda = 2$, $4$, $8$ and $16$ on the forms of the sixth analysis. As controls on synthetic writers, one
chunk holding the whole log reproduces the costs of A2 on windows when its start is fixed and on full texts when both
ends are fixed.

\paragraph{The strongest attacker with the variants.}
For each form and $N$, the strongest attacker with the variants is the maximum over the attackers of the main
computation and the variants of the sixth and seventh analyses, selected again in every resample. Selecting on the
evaluated logs favours the attacker (\S\ref{sec:threat-model:bias}). Thirteen of the fourteen pairs of form and $N$
that a variant won were won with $\lambda = 2$, the lowest value of the grid, and the other with $\lambda = 16$
(Table~\ref{tab:app:stronger}); a lighter penalty might match more, and was not tried. The process measures do not
change, since the forms do not.

\paragraph{Tables.}
% The computation of R-DP (run_main.py --forms dp) ran R3 and R-DP without R0, so the summaries of the process measures
% relative to R0 were empty in exploratory_dp.json (notes.md, D5). They were recomputed from the caches with R0 of the
% main computation, with the same sample and resamples (exploratory_dp_utility_r0.json); the values of R3 equal those
% of main.json.
Tables~\ref{tab:app:unicity} to~\ref{tab:app:band} give the first five analyses, and
Tables~\ref{tab:app:lambda-1000} to~\ref{tab:app:stronger} the last two.

\begin{table*}
  \caption{Exploratory: uniqueness of word-length sequences among $N$ candidates, 500 writers. Unique: the percentage
  of runs of $k$ consecutive words whose sequence of word lengths no decoy contains, averaged per writer and then over
  writers, with its interval. $1/m$: the accuracy of choosing uniformly among the $m$ candidates that contain the
  sequence. All: unique at $N = 1{,}000$ for all writers after exclusion. $k^{*}$: the reference lines $\log_2 N / h_j$ for
  $j = 0$, $1$, $2$ preceding words, with $h_j$ the plug-in entropy rate of word lengths over all writers after exclusion.}
  \label{tab:app:unicity}
  \small
  \begin{tabular}{rllllllr}
    \toprule
    & \multicolumn{2}{l}{$N = 10$} & \multicolumn{2}{l}{$N = 100$} & \multicolumn{2}{l}{$N = 1{,}000$} & All \\
    $k$ & Unique & $1/m$ & Unique & $1/m$ & Unique & $1/m$ & $1{,}000$ \\
    \midrule
    \tabinput{sections/generated/explore-unicity.tex}
    \bottomrule
  \end{tabular}
\end{table*}

\begin{table*}
  \caption{Exploratory: R-DP($\varepsilon$), with $\varepsilon$ the budget per segment count and per interval, not per
  writer, and R3, 500 writers. Full texts: the accuracy of the strongest attacker, A1$'$ on every form and $N$, and the
  contrast R3~$-$~R-DP in points at $N = 1{,}000$. Windows: the accuracy of the strongest attacker, A2, at
  $N = 1{,}000$. Area: the open-world area. The values of R3 are those of the main computation.}
  \label{tab:app:dp}
  \small
  \begin{tabular}{lrrlllllll}
    \toprule
    & \multicolumn{4}{l}{Full texts} & \multicolumn{4}{l}{Windows, $N = 1{,}000$} & \\
    Form & $N = 10$ & $100$ & $1{,}000$ & R3~$-$~form & $k = 5$ & $10$ & $20$ & $50$ & Area \\
    \midrule
    \tabinput{sections/generated/explore-dp.tex}
    \bottomrule
  \end{tabular}
\end{table*}

\begin{table*}
  \caption{Exploratory: process measures on R3 and R-DP($\varepsilon$), 500 writers, mean over the four prompts. Top:
  the Spearman correlation with the value in R0. Bottom: the Spearman correlation with the human score. Every measure
  on R-DP is approximated.}
  \label{tab:app:dp-score}
  \small
  \begin{tabular}{lrrrrr}
    \toprule
    Measure & R3 & R-DP(4) & R-DP(2) & R-DP(1) & R-DP(0.5) \\
    \midrule
    \multicolumn{6}{l}{Correlation with R0} \\
    \tabinput{sections/generated/explore-dp-r0.tex}
    \midrule
    \multicolumn{6}{l}{Correlation with the human score} \\
    \tabinput{sections/generated/explore-dp-score.tex}
    \bottomrule
  \end{tabular}
\end{table*}

\begin{table*}
  \caption{Exploratory: the upper tail of accuracy per writer, full texts, 500 writers: the 95th and 99th percentiles
  and the share of writers whose true text was chosen alone in every draw.}
  \label{tab:app:tail}
  \small
    \begin{tabular}{lrllll}
    \toprule
    Form & $N$ & Attacker & 95th & 99th & Every draw \\
    \midrule
    \tabinput{sections/generated/explore-tail.tex}
    \bottomrule
  \end{tabular}
\end{table*}

\begin{table*}
  \caption{Exploratory: matches above a margin, full texts, 500 writers. For thresholds at the 0th, 20th, 50th, 80th
  and 90th percentiles of the margin between the second-lowest and the lowest cost, the share of logs and draws at or
  above the threshold and the accuracy among them (share/accuracy). Costs of different attackers are not comparable.}
  \label{tab:app:margin}
  \small
  \begin{tabular}{lrllllll}
    \toprule
    Form & $N$ & Attacker & 0th & 20th & 50th & 80th & 90th \\
    \midrule
    \tabinput{sections/generated/explore-margin.tex}
    \bottomrule
  \end{tabular}
\end{table*}

\begin{table*}
  \caption{Exploratory: accuracy by score band (thirds of each prompt's sampled writers, lowest first), full texts,
  and the difference between the highest and the lowest third.}
  \label{tab:app:band}
  \small
    \begin{tabular}{lrlrrrl}
    \toprule
    Form & $N$ & Att. & Low & Mid & High & High $-$ low \\
    \midrule
    \tabinput{sections/generated/explore-band.tex}
    \bottomrule
  \end{tabular}
\end{table*}

\begin{table*}
  \caption{Exploratory: accuracy (\%) of every variant of the sixth and seventh analyses, full texts, $N = 1{,}000$, with
  intervals. A2 and A1$'$ with $\lambda = 8$ are the attackers of the main computation; A2-scan aligns chunks of 20
  segments anywhere in the text.}
  \label{tab:app:lambda-1000}
  \small
  \begin{tabular}{llllll}
    \toprule
    Form & Attacker & $\lambda = 2$ & $\lambda = 4$ & $\lambda = 8$ & $\lambda = 16$ \\
    \midrule
    \tabinput{sections/generated/explore-lambda-scan-1000.tex}
    \bottomrule
  \end{tabular}
\end{table*}

\begin{table*}
  \caption{Exploratory: accuracy (\%) of every variant of the sixth and seventh analyses, full texts, $N = 10$, with
  intervals. A2 and A1$'$ with $\lambda = 8$ are the attackers of the main computation; A2-scan aligns chunks of 20
  segments anywhere in the text.}
  \label{tab:app:lambda-10}
  \small
  \begin{tabular}{llllll}
    \toprule
    Form & Attacker & $\lambda = 2$ & $\lambda = 4$ & $\lambda = 8$ & $\lambda = 16$ \\
    \midrule
    \tabinput{sections/generated/explore-lambda-scan-10.tex}
    \bottomrule
  \end{tabular}
\end{table*}

\begin{table*}
  \caption{Exploratory: the strongest attacker of the main computation (Main) and the strongest attacker with the
  variants of the sixth and seventh analyses (With variants, selected again in every resample, with its interval and
  the attacker chosen, with the percentage of resamples in parentheses when below 100), full texts, accuracy in
  percent. Difference: with variants minus main at $N = 1{,}000$, with its interval.}
  \label{tab:app:stronger}
  \small
  \begin{tabular}{lllllllr}
    \toprule
    & \multicolumn{3}{l}{$N = 10$} & \multicolumn{4}{l}{$N = 1{,}000$} \\
    Form & Main & With variants & Attacker & Main & With variants & Attacker & Difference \\
    \midrule
    \tabinput{sections/generated/explore-strongest.tex}
    \bottomrule
  \end{tabular}
\end{table*}
